\chapter{Región de Interés (ROI)}
\label{ch:roi}
% ==============================================================

\section{Criterios de delimitación}

Para optimizar la segmentación con SAM y reducir ruido en la extracción de línea de costa,
cada cámara tiene una ROI que excluye:

\begin{itemize}[leftmargin=1.5em]
  \item \textbf{Cielo}: píxeles por encima del horizonte.
  \item \textbf{Infraestructura}: paseos marítimos, espigones, edificios fuera de la zona de interés.
  \item \textbf{Mar profundo}: zona sin rotura de onda donde la línea de costa no es visible.
\end{itemize}

\section{Coordenadas por cámara}

Las ROIs se almacenan en \texttt{calibration/roi\_config.json} como un diccionario por
cámara con claves \texttt{x\_min}, \texttt{y\_min}, \texttt{x\_max}, \texttt{y\_max} (píxeles,
sobre la resolución nativa de captura de cada cámara — \textbf{4608$\times$2682}, no
1920$\times$1080). La tabla~\ref{tab:roi} recoge los valores de producción actualmente
configurados.

\begin{table}[H]
\centering
\caption{ROIs de producción por cámara (calibration/roi\_config.json)}
\label{tab:roi}
\begin{tabular}{C{2cm} C{4.5cm} L{6cm}}
\toprule
\textbf{Cámara} & \textbf{ROI $(x_{min}, y_{min}, x_{max}, y_{max})$ px} & \textbf{Justificación} \\
\midrule
CAM~1 & [0, 993, 4608, 2682]  & Excluye horizonte y cielo \\
CAM~2 & [0, 869, 4608, 2682]  & Horizonte más bajo \\
CAM~3 & [0, 1118, 4608, 2682] & Ángulo con más cielo visible \\
CAM~4 & [0, 993, 4608, 2682]  & Excluye horizonte y cielo \\
CAM~5 & [0, 944, 4608, 2682]  & Excluye horizonte y cielo \\
CAM~6 & [0, 1043, 4608, 2682] & Excluye horizonte y cielo \\
\bottomrule
\end{tabular}
\end{table}

\begin{infobox}[¿De dónde sale un valor por defecto si falta la ROI de una cámara?]
Si \texttt{roi\_config.json} no tiene entrada para una cámara, \texttt{analyze\_roi()} calcula
una por defecto recortando el \textbf{40\% superior} de la imagen (\texttt{y\_min = 0.4 ×
alto}) — una aproximación razonable al horizonte mientras no se afine a mano, no un valor
definitivo. Los seis valores de la tabla ya están afinados individualmente y sustituyen a
ese valor por defecto.
\end{infobox}

% ==============================================================
